\subsection{Continuous laws and their averages}\label{subsec:measures-laws-and-averages}
\paraid{p-e99b54a5841a}
We now select probability laws on the fibers of the forgetful map
and average the invariant measures with respect to these laws.

\paraid{p-f921de5bf828}
The averaging map in \Cref{lem:continuous-barycenter} is the standard
barycenter map on the space of probability measures.
Its construction and weak continuity appear in
\cite[proof of Theorem~8.10.5, p.~218]{Bogachev2007II}.
We include a proof using the Riesz--Markov--Kakutani theorem
(\Cref{thm:riesz-markov-kakutani}),
together with the full-support and invariance properties needed below.
We then prove continuity as the compact metric space varies by comparing
integrals in a common compact ambient space.

\begin{proposition}\label{lem:continuous-barycenter}
\paraid{p-ee76895903f0}
Let
$\MetricSpace{\BaseCarrier}{\MetricSymbol}$
be a nonempty compact metric space.
There exists a unique map
\[
 \operatorname{bar}_\BaseCarrier\colon
 \ProbabilityMeasures(\ProbabilityMeasures(\BaseCarrier))
 \to\ProbabilityMeasures(\BaseCarrier)
\]
such that for every
$\FirstProbabilityLaw\in\ProbabilityMeasures(\ProbabilityMeasures(\BaseCarrier))$
and every
$\TestFunction\in\ContinuousFunctions(\BaseCarrier)$,
\begin{equation}\label{eq:continuous-barycenter}
 \int_\BaseCarrier\TestFunction\,
       \IntegrationDifferential\operatorname{bar}_\BaseCarrier(\FirstProbabilityLaw)
 =\int_{\ProbabilityMeasures(\BaseCarrier)}
       \left(\int_\BaseCarrier\TestFunction\,\IntegrationDifferential\FirstMeasure\right)
       \,\IntegrationDifferential\FirstProbabilityLaw(\FirstMeasure).
\end{equation}
This map is continuous for the weak topologies.
If
$\FirstProbabilityLaw$
is concentrated on a Borel set of full-support measures,
then
$\operatorname{bar}_\BaseCarrier(\FirstProbabilityLaw)$
has full support.
For each
$\GroupElement\in\Isom\MetricSpace{\BaseCarrier}{\MetricSymbol}$,
assume that
\begin{equation}\label{eq:barycenter-invariant-law}
 \FirstProbabilityLaw\left(
 \left\{\FirstMeasure\in\ProbabilityMeasures(\BaseCarrier)
 \,\middle|\,
 \GroupElement\Pushforward\FirstMeasure=\FirstMeasure\right\}
 \right)=1.
\end{equation}
Then
\[
 \GroupElement\Pushforward\operatorname{bar}_\BaseCarrier(\FirstProbabilityLaw)
 =\operatorname{bar}_\BaseCarrier(\FirstProbabilityLaw).
\]
\end{proposition}

\begin{proof}
\paraid{p-11561faad095}
For
$\TestFunction\in\ContinuousFunctions(\BaseCarrier)$,
define
$\widehat\TestFunction\colon\ProbabilityMeasures(\BaseCarrier)\to\RealNumbers$
as follows.
For every
$\FirstMeasure\in\ProbabilityMeasures(\BaseCarrier)$,
set
\[
 \widehat\TestFunction(\FirstMeasure)
 =\int_\BaseCarrier\TestFunction\,\IntegrationDifferential\FirstMeasure.
\]
This function is continuous and bounded in absolute value by
$\norm{\TestFunction}_\infty$.
Fix
$\FirstProbabilityLaw\in\ProbabilityMeasures(\ProbabilityMeasures(\BaseCarrier))$.
Define the functional
$\RepresentingFunctional\colon\ContinuousFunctions(\BaseCarrier)\to\RealNumbers$
by integrating these functions,
\[
 \RepresentingFunctional(\TestFunction)
 =\int_{\ProbabilityMeasures(\BaseCarrier)}\widehat\TestFunction
       \,\IntegrationDifferential\FirstProbabilityLaw.
\]
This functional is positive and linear,
and
$\RepresentingFunctional(1)=1$.
\Cref{thm:riesz-markov-kakutani} therefore yields a unique representing
probability satisfying \eqref{eq:continuous-barycenter}.
If
$\FirstProbabilityLaw_\SequenceIndex\to\FirstProbabilityLaw$
weakly in
$\ProbabilityMeasures(\ProbabilityMeasures(\BaseCarrier))$,
then for every
$\TestFunction\in\ContinuousFunctions(\BaseCarrier)$,
\[
 \int_{\ProbabilityMeasures(\BaseCarrier)}\widehat\TestFunction
       \,\IntegrationDifferential\FirstProbabilityLaw_\SequenceIndex
 \longrightarrow
 \int_{\ProbabilityMeasures(\BaseCarrier)}\widehat\TestFunction
       \,\IntegrationDifferential\FirstProbabilityLaw.
\]
Thus the barycenters converge weakly.
Since both probability spaces are metrizable in their weak topologies,
this proves continuity.

\paraid{p-055551975ef9}
Assume that
$\FirstProbabilityLaw$
is concentrated on a Borel set of full-support measures.
For a nonempty open subset
$\OpenSubset\subset\BaseCarrier$,
choose a continuous function
$\TestFunction\colon\BaseCarrier\to[0,1]$
that is nonzero and vanishes outside
$\OpenSubset$.
For every
$\FirstMeasure\in\ProbabilityMeasures(\BaseCarrier)$
with
$\supp(\FirstMeasure)=\BaseCarrier$,
we have
$\widehat\TestFunction(\FirstMeasure)>0$.
By the hypothesis and \eqref{eq:continuous-barycenter},
\[
 \operatorname{bar}_\BaseCarrier(\FirstProbabilityLaw)(\OpenSubset)
 \geq\int_{\ProbabilityMeasures(\BaseCarrier)}
       \widehat\TestFunction(\FirstMeasure)
       \,\IntegrationDifferential\FirstProbabilityLaw(\FirstMeasure)>0.
\]
This proves full support.

\paraid{p-5cc5a2ea18dd}
Fix an isometry
$\GroupElement\in\Isom\MetricSpace{\BaseCarrier}{\MetricSymbol}$
for which \eqref{eq:barycenter-invariant-law} holds.
For every
$\TestFunction\in\ContinuousFunctions(\BaseCarrier)$
and every
$\FirstMeasure\in\ProbabilityMeasures(\BaseCarrier)$
with
$\GroupElement\Pushforward\FirstMeasure=\FirstMeasure$,
we have
\begin{equation}\label{eq:fiber-invariant-integrals}
 \int_\BaseCarrier\TestFunction\circ\GroupElement\,
       \IntegrationDifferential\FirstMeasure
 =\int_\BaseCarrier\TestFunction\,
       \IntegrationDifferential\FirstMeasure.
\end{equation}
By \eqref{eq:barycenter-invariant-law},
integrating \eqref{eq:fiber-invariant-integrals} against
$\FirstProbabilityLaw$
and using \eqref{eq:continuous-barycenter} proves invariance of
$\operatorname{bar}_\BaseCarrier(\FirstProbabilityLaw)$
under
$\GroupElement$.
\end{proof}

\begin{theorem}\label{thm:invariant-measures}
\paraid{p-b7a56605e33d}
For every nonempty compact metric space
$\MetricSpace{\BaseCarrier}{\MetricSymbol}$,
there exists a probability measure
$\SelectedMeasure{\BaseCarrier}{\MetricSymbol}\in\ProbabilityMeasures(\BaseCarrier)$.
These measures can be chosen simultaneously to satisfy the three properties below.
\begin{enumerate}[label=\textup{(M\arabic*)},ref=\textup{(M\arabic*)}]
\item
\paraid{p-09630bcb2594}
For every
$\MetricSpace{\BaseCarrier}{\MetricSymbol}$,
we have
$\supp(\SelectedMeasure{\BaseCarrier}{\MetricSymbol})=\BaseCarrier$.
\item
\paraid{p-1eec8fcfc366}
For every pair of nonempty compact metric spaces
$\MetricSpace{\BaseCarrier}{\MetricSymbol}$
and
$\MetricSpace{\ComparisonCarrier }{\ComparisonMetric }$
and every isometry
$\IdentifyingIsometry \colon\MetricSpace{\BaseCarrier }{\MetricSymbol }\to\MetricSpace{\ComparisonCarrier }{\ComparisonMetric }$,
we have
\[
 \IdentifyingIsometry \Pushforward\SelectedMeasure{\BaseCarrier}{\MetricSymbol}=\SelectedMeasure{\ComparisonCarrier}{\ComparisonMetric}.
\]
\item\label{item:invariant-measures-continuity}
\paraid{p-4f2c8f1817cc}
Let
$\MetricSpace{\BaseCarrier _\SequenceIndex }{\MetricSymbol _\SequenceIndex }$
and
$\MetricSpace{\BaseCarrier }{\MetricSymbol }$
be nonempty compact metric spaces.
Let
$\MetricSpace{\AmbientCarrier}{\AmbientMetric}$
be any compact metric space and let
$\FirstEmbedding _\SequenceIndex \colon\MetricSpace{\BaseCarrier _\SequenceIndex }{\MetricSymbol _\SequenceIndex }\to\MetricSpace{\AmbientCarrier }{\AmbientMetric }$
and
$\FirstEmbedding \colon\MetricSpace{\BaseCarrier }{\MetricSymbol }\to\MetricSpace{\AmbientCarrier }{\AmbientMetric }$
be isometric embeddings.
Then
\[
 \HausdorffDistance{\AmbientMetric }(\FirstEmbedding _\SequenceIndex (\BaseCarrier _\SequenceIndex ),\FirstEmbedding (\BaseCarrier ))\to0
 \quad\Longrightarrow\quad
 (\FirstEmbedding _\SequenceIndex )\Pushforward\SelectedMeasure{\BaseCarrier _\SequenceIndex}{\MetricSymbol _\SequenceIndex}\to\FirstEmbedding \Pushforward\SelectedMeasure{\BaseCarrier}{\MetricSymbol}
 \quad\text{weakly in }\ProbabilityMeasures(\AmbientCarrier ).
\]
\end{enumerate}
\end{theorem}

\begin{proof}
\paraid{p-941b13a492c1}
By \Cref{prop:invariant-measured-polish,prop:open-invariant-projection},
the map
$\MeasureProjection\colon\InvariantMeasuredSpaces\to\GHSpace$
is an open continuous surjection between Polish spaces.
By Valov's theorem (\Cref{thm:valov}),
we obtain a continuous map
$\LawSection\colon\GHSpace\to\ProbabilityMeasures(\InvariantMeasuredSpaces)$
such that for every
$\MetricSpace{\BaseCarrier}{\MetricSymbol}\in\GHSpace$
we have
$\MeasureProjection\Pushforward\LawSection\MetricSpace{\BaseCarrier }{\MetricSymbol }=\DiracProbability _{\MetricSpace{\BaseCarrier }{\MetricSymbol }}$.
Fix a representative
$\MetricSpace{\BaseCarrier }{\MetricSymbol }$
and set
\[
 \MeasureFiber{\BaseCarrier }=\MeasureProjection^{-1}(\MetricSpace{\BaseCarrier }{\MetricSymbol }).
\]
This is a closed Polish subspace of
$\InvariantMeasuredSpaces$,
and
$\LawSection\MetricSpace{\BaseCarrier }{\MetricSymbol }(\MeasureFiber{\BaseCarrier })=1$.
We construct a map
$\FiberMeasureMap{\BaseCarrier}\colon\MeasureFiber{\BaseCarrier}\to\ProbabilityMeasures(\BaseCarrier)$
on the fixed carrier as follows.
For each class
$\FiberPoint\in\MeasureFiber{\BaseCarrier}$,
choose a representative
$\MeasuredSpace{\ComparisonCarrier}{\ComparisonMetric}{\SecondMeasure}$
of $\FiberPoint$ and an isometry
$\IdentifyingIsometry \colon\MetricSpace{\ComparisonCarrier }{\ComparisonMetric }\to\MetricSpace{\BaseCarrier }{\MetricSymbol }$.
We denote the pushforward measure on this fixed representative by
\[
 \FiberMeasure{\BaseCarrier}{\FiberPoint}=\IdentifyingIsometry\Pushforward\SecondMeasure.
\]
For every
$\BorelSubset\in\BorelSets(\BaseCarrier)$,
this means that
\[
 \FiberMeasure{\BaseCarrier}{\FiberPoint}(\BorelSubset)
 =\SecondMeasure(\IdentifyingIsometry^{-1}(\BorelSubset)).
\]
Set
$\FiberMeasureMap{\BaseCarrier}(\FiberPoint)=\FiberMeasure{\BaseCarrier}{\FiberPoint}$.
If
$\SecondIdentifyingIsometry \colon\MetricSpace{\ComparisonCarrier }{\ComparisonMetric }\to\MetricSpace{\BaseCarrier }{\MetricSymbol }$
is another isometry,
then
$\SecondIdentifyingIsometry ^{-1}\IdentifyingIsometry \in\Isom\MetricSpace{\ComparisonCarrier }{\ComparisonMetric }$
and
\[
 \IdentifyingIsometry \Pushforward\SecondMeasure
 =\SecondIdentifyingIsometry \Pushforward\bigl((\SecondIdentifyingIsometry ^{-1}\IdentifyingIsometry )\Pushforward\SecondMeasure \bigr)
 =\SecondIdentifyingIsometry \Pushforward\SecondMeasure .
\]
Thus the measure is independent of the chosen isometry from this representative.
Now let $(Y_0,e_0,\nu_0)$ be another representative of $\FiberPoint$.
Choose a measure-preserving isometry
$h\colon\MetricSpace{\ComparisonCarrier}{\ComparisonMetric}\to(Y_0,e_0)$,
so that $h_*\SecondMeasure=\nu_0$,
and an isometry $k\colon(Y_0,e_0)\to\MetricSpace{\BaseCarrier}{\MetricSymbol}$.
Independence of the transport isometry on $\ComparisonCarrier$ implies
\[
 k_*\nu_0=(k\circ h)_*\SecondMeasure
 =\IdentifyingIsometry_*\SecondMeasure.
\]
Hence the measure is also independent of the representative of $\FiberPoint$.

\paraid{p-218d0d8ad567}
We first prove continuity of this assignment of measures under the hypotheses of \ref{item:invariant-measures-continuity}.
Fix a sequence
$\{\MetricSpace{\BaseCarrier_\SequenceIndex}{\MetricSymbol_\SequenceIndex}\}_{\SequenceIndex\in\NonnegativeIntegers}$,
a nonempty compact metric space
$\MetricSpace{\BaseCarrier}{\MetricSymbol}$,
a compact metric space
$\MetricSpace{\AmbientCarrier}{\AmbientMetric}$,
and prescribed isometric embeddings
$\FirstEmbedding_\SequenceIndex$
and
$\FirstEmbedding$
as in \ref{item:invariant-measures-continuity}, with
\[
 \HausdorffDistance{\AmbientMetric}(\FirstEmbedding_\SequenceIndex(\BaseCarrier_\SequenceIndex),\FirstEmbedding(\BaseCarrier))\to0.
\]
Identify
$\BaseCarrier_\SequenceIndex$
and
$\BaseCarrier$
with their images under the isometric embeddings into
$\AmbientCarrier$.
If
$\FiberPoint _\SequenceIndex \in\MeasureFiber{\BaseCarrier _\SequenceIndex }$
and
$\FiberPoint \in\MeasureFiber{\BaseCarrier }$
satisfy
$\FiberPoint _\SequenceIndex \to\FiberPoint $
in
$\InvariantMeasuredSpaces$,
then
\begin{equation}\label{eq:fiber-pullbacks}
 \FiberMeasure{\BaseCarrier _\SequenceIndex}{\FiberPoint _\SequenceIndex}\longrightarrow\FiberMeasure{\BaseCarrier}{\FiberPoint}
 \quad\text{weakly in }\ProbabilityMeasures(\AmbientCarrier ).
\end{equation}
Indeed,
let
$\SecondMeasure$
be the limit of any weakly convergent subsequence
\[
 \{\FiberMeasure{\BaseCarrier_{n_k}}{\FiberPoint_{n_k}}\}_{k\in\NonnegativeIntegers}.
\]
Such subsequences exist by compactness of
$\ProbabilityMeasures(\AmbientCarrier)$.
The limit is concentrated on
$\BaseCarrier $,
by \Cref{lem:limit-support}.
By the definition of
$\GHPDistance$
and weak convergence on
$\AmbientCarrier $,
we obtain
\[
 \FiberPoint _{n_k} \longrightarrow\MeasuredSpace{\BaseCarrier }{\MetricSymbol }{\SecondMeasure }.
\]
By uniqueness of limits,
$\MeasuredSpace{\BaseCarrier }{\MetricSymbol }{\SecondMeasure }=\FiberPoint $.
Thus there exists an isometry
$\GroupElement\colon\BaseCarrier\to\BaseCarrier$
such that
$\SecondMeasure=\GroupElement\Pushforward\FiberMeasure{\BaseCarrier}{\FiberPoint}$.
The measure
$\FiberMeasure{\BaseCarrier}{\FiberPoint}$
is invariant,
so
\begin{equation}\label{eq:fiber-measure-limit}
 \SecondMeasure=\FiberMeasure{\BaseCarrier}{\FiberPoint}.
\end{equation}
Compactness of
$\ProbabilityMeasures(\AmbientCarrier)$
and uniqueness of the subsequential limit in
\eqref{eq:fiber-measure-limit} imply \eqref{eq:fiber-pullbacks}.
With the space
$\BaseCarrier $
fixed,
this also proves continuity of
$\FiberMeasureMap{\BaseCarrier }$.

\paraid{p-39f42cc41be0}
Push the law on the fiber forward by
$\FiberMeasureMap{\BaseCarrier}$
and define
\[
 \FirstProbabilityLaw_\BaseCarrier
 =(\FiberMeasureMap{\BaseCarrier})\Pushforward
   \left(\LawSection\MetricSpace{\BaseCarrier}{\MetricSymbol}
                 |_{\MeasureFiber{\BaseCarrier}}\right).
\]
Using the barycenter in \Cref{lem:continuous-barycenter},
we define
\[
 \SelectedMeasure{\BaseCarrier}{\MetricSymbol}
 =\operatorname{bar}_\BaseCarrier(\FirstProbabilityLaw_\BaseCarrier).
\]
Define
$\RepresentingFunctional_\BaseCarrier\colon\ContinuousFunctions(\BaseCarrier)\to\RealNumbers$
as follows.
For every
$\TestFunction\in\ContinuousFunctions(\BaseCarrier)$,
set
\begin{equation}\label{eq:average-functional}
 \RepresentingFunctional_\BaseCarrier(\TestFunction)
 =\int_{\MeasureFiber{\BaseCarrier}}
       \left(\int_\BaseCarrier\TestFunction\,
       \IntegrationDifferential\FiberMeasure{\BaseCarrier}{\FiberPoint}\right)
       \IntegrationDifferential\LawSection
       \MetricSpace{\BaseCarrier}{\MetricSymbol}(\FiberPoint).
\end{equation}
Substituting
$\FirstProbabilityLaw_\BaseCarrier$
for
$\FirstProbabilityLaw$
in the identity \eqref{eq:continuous-barycenter} gives
\begin{equation}\label{eq:measure-barycenter}
 \int_\BaseCarrier\TestFunction\,
       \IntegrationDifferential\SelectedMeasure{\BaseCarrier}{\MetricSymbol}
 =\RepresentingFunctional_\BaseCarrier(\TestFunction).
\end{equation}
Figure~\ref{fig:measure-average} shows the two levels of probability measures.
\begin{figure}[H]
\centering
\begin{tikzpicture}[>=Stealth,node distance=9mm,font=\small]
\node (law) {$\LawSection\MetricSpace{\BaseCarrier}{\MetricSymbol}
 \in\ProbabilityMeasures(\MeasureFiber{\BaseCarrier})$};
\node[below=of law] (transport) {$
 (\FiberMeasureMap{\BaseCarrier})_*\LawSection\MetricSpace{\BaseCarrier}{\MetricSymbol}
 \in\ProbabilityMeasures(\ProbabilityMeasures(\BaseCarrier))$};
\node[below=of transport] (average) {$
 \SelectedMeasure{\BaseCarrier}{\MetricSymbol}\in\ProbabilityMeasures(\BaseCarrier)$};
\draw[->] (law) -- node[right]{transport by $\FiberMeasureMap{\BaseCarrier}$} (transport);
\draw[->] (transport) -- node[right]{average} (average);
\end{tikzpicture}
\caption{For a fixed representative
$\MetricSpace{\BaseCarrier}{\MetricSymbol}$,
the law
$\LawSection\MetricSpace{\BaseCarrier}{\MetricSymbol}$
on the fiber
$\MeasureFiber{\BaseCarrier}$
is transported to a law on
$\ProbabilityMeasures(\BaseCarrier)$.
The averaged measure
$\SelectedMeasure{\BaseCarrier}{\MetricSymbol}$
is defined by \eqref{eq:average-functional} and \eqref{eq:measure-barycenter}.}
\label{fig:measure-average}
\end{figure}


\paraid{p-4bf68c702f9c}
For a compact metric space
$\MetricSpace{\BaseCarrier}{\MetricSymbol}$,
we denote by
$\FullSupportProbabilities{\BaseCarrier}$
the set of full-support probabilities on
$\BaseCarrier$,
defined by
\[
 \FullSupportProbabilities{\BaseCarrier}
 =\{\FirstMeasure\in\ProbabilityMeasures(\BaseCarrier)
       \mid\supp\FirstMeasure=\BaseCarrier\}.
\]
For each
$\GroupElement\in\Isom\MetricSpace{\BaseCarrier}{\MetricSymbol}$,
we denote by
$\InvariantProbabilities{\BaseCarrier}{\GroupElement}$
the set of probabilities invariant under
$\GroupElement$,
defined by
\[
 \InvariantProbabilities{\BaseCarrier}{\GroupElement}
 =\{\FirstMeasure\in\ProbabilityMeasures(\BaseCarrier)
       \mid\GroupElement\Pushforward\FirstMeasure=\FirstMeasure\}.
\]
The law
$\FirstProbabilityLaw_\BaseCarrier$
is concentrated on
\[
 \FullSupportProbabilities{\BaseCarrier}
 \cap\bigcap_{\GroupElement\in\Isom\MetricSpace{\BaseCarrier}{\MetricSymbol}}
       \InvariantProbabilities{\BaseCarrier}{\GroupElement}.
\]
Choose a countable base $\CountableBasis$ for
$\BaseCarrier$.
A probability
$\FirstMeasure\in\ProbabilityMeasures(\BaseCarrier)$
has full support if and only if
$\FirstMeasure(\OpenSubset)>0$
for every nonempty
$\OpenSubset\in\CountableBasis$.
The weak topology on
$\ProbabilityMeasures(\BaseCarrier)$
is metrizable.
For each nonempty
$\OpenSubset\in\CountableBasis$
and each sequence
$\{\FirstMeasure_\SequenceIndex\}_{\SequenceIndex\in\NonnegativeIntegers}$
that converges weakly to
$\FirstMeasure\in\ProbabilityMeasures(\BaseCarrier)$,
\Cref{thm:portmanteau}
applied to the open set
$\OpenSubset$
gives
\[
 \FirstMeasure(\OpenSubset)
 \leq\liminf_{\SequenceIndex\to\infty}
          \FirstMeasure_\SequenceIndex(\OpenSubset).
\]
The sequential criterion for lower semicontinuity in a metrizable space now
shows that the map
\[
 \ProbabilityMeasures(\BaseCarrier)\to[0,1],
 \qquad\FirstMeasure\mapsto\FirstMeasure(\OpenSubset)
\]
is lower semicontinuous for the weak topology.
The strict superlevel set at $0$ is open, so
\[
 \{\FirstMeasure\in\ProbabilityMeasures(\BaseCarrier)
      \mid \FirstMeasure(\OpenSubset)>0\}
\]
is open in
$\ProbabilityMeasures(\BaseCarrier)$.
The basis criterion yields
\[
 \FullSupportProbabilities{\BaseCarrier}
 =\bigcap_{\substack{\OpenSubset\in\CountableBasis\\
                      \OpenSubset\ne\emptyset}}
       \{\FirstMeasure\in\ProbabilityMeasures(\BaseCarrier)
           \mid\FirstMeasure(\OpenSubset)>0\}.
\]
Since the basis is countable and each set in this intersection is open,
$\FullSupportProbabilities{\BaseCarrier}$
is a countable intersection of open sets, and in particular is Borel.
For each
$\GroupElement\in\Isom\MetricSpace{\BaseCarrier}{\MetricSymbol}$,
the set
$\InvariantProbabilities{\BaseCarrier}{\GroupElement}$
is closed in
$\ProbabilityMeasures(\BaseCarrier)$,
as we now verify.
For every
$\FirstMeasure\in\ProbabilityMeasures(\BaseCarrier)$
and every
$\TestFunction\in\ContinuousFunctions(\BaseCarrier)$,
the pushforward satisfies
\[
 \int_\BaseCarrier\TestFunction\,
       \IntegrationDifferential(\GroupElement\Pushforward\FirstMeasure)
 =\int_\BaseCarrier(\TestFunction\circ\GroupElement)\,
       \IntegrationDifferential\FirstMeasure.
\]
Since
$\GroupElement$
is continuous,
$\TestFunction\circ\GroupElement$
is continuous.
Let
$\{\FirstMeasure_\SequenceIndex\}_{\SequenceIndex\in\NonnegativeIntegers}$
be a sequence in
$\ProbabilityMeasures(\BaseCarrier)$
that converges weakly to
$\FirstMeasure$.
For every
$\TestFunction\in\ContinuousFunctions(\BaseCarrier)$,
the integral identity gives
\[
 \int_\BaseCarrier\TestFunction\,
       \IntegrationDifferential(\GroupElement\Pushforward\FirstMeasure_\SequenceIndex)
 \longrightarrow
 \int_\BaseCarrier\TestFunction\,
       \IntegrationDifferential(\GroupElement\Pushforward\FirstMeasure).
\]
Hence the pushforward map
\[
 \GroupElement\Pushforward\colon
 \ProbabilityMeasures(\BaseCarrier)\to\ProbabilityMeasures(\BaseCarrier)
\]
is continuous for the weak topology.
Now let
$\{\FirstMeasure_\SequenceIndex\}_{\SequenceIndex\in\NonnegativeIntegers}$
be a sequence in
$\InvariantProbabilities{\BaseCarrier}{\GroupElement}$
that converges weakly to
$\FirstMeasure\in\ProbabilityMeasures(\BaseCarrier)$.
Continuity and invariance give
\[
 \GroupElement\Pushforward\FirstMeasure_\SequenceIndex
 =\FirstMeasure_\SequenceIndex\longrightarrow\FirstMeasure,
 \qquad
 \GroupElement\Pushforward\FirstMeasure_\SequenceIndex
 \longrightarrow\GroupElement\Pushforward\FirstMeasure.
\]
Uniqueness of weak limits gives
$\GroupElement\Pushforward\FirstMeasure=\FirstMeasure$.
Thus
$\InvariantProbabilities{\BaseCarrier}{\GroupElement}$
is sequentially closed.
Since the weak topology is metrizable,
$\InvariantProbabilities{\BaseCarrier}{\GroupElement}$
is closed.
\Cref{lem:continuous-barycenter} therefore implies that
$\SelectedMeasure{\BaseCarrier}{\MetricSymbol}$
is invariant and has full support.
Finally,
for an isometry
$\IdentifyingIsometry \colon\MetricSpace{\BaseCarrier }{\MetricSymbol }\to\MetricSpace{\ComparisonCarrier }{\ComparisonMetric }$,
we have the equality of isometry classes
\[
 \MetricSpace{\BaseCarrier}{\MetricSymbol}
 =\MetricSpace{\ComparisonCarrier}{\ComparisonMetric}
 \quad\text{in }\GHSpace.
\]
Thus the fibers satisfy
\[
 \MeasureFiber{\BaseCarrier}=\MeasureFiber{\ComparisonCarrier}
 \quad\text{as subsets of }\InvariantMeasuredSpaces,
\]
and,
since the map
$\LawSection$
is defined on GH isometry classes,
\[
 \LawSection\MetricSpace{\BaseCarrier}{\MetricSymbol}
 =\LawSection\MetricSpace{\ComparisonCarrier}{\ComparisonMetric}.
\]
Fix the representatives
$\MetricSpace{\BaseCarrier}{\MetricSymbol}$
and
$\MetricSpace{\ComparisonCarrier}{\ComparisonMetric}$
and the isometry
$\IdentifyingIsometry$
between them.
For every
$\FiberPoint\in\MeasureFiber{\BaseCarrier}$,
the definition by transport implies
\[
 \FiberMeasure{\ComparisonCarrier}{\FiberPoint}
 =\IdentifyingIsometry\Pushforward\FiberMeasure{\BaseCarrier}{\FiberPoint}.
\]
For every
$\TestFunction\in\ContinuousFunctions(\ComparisonCarrier)$,
apply \eqref{eq:average-functional} and \eqref{eq:measure-barycenter} on
$\BaseCarrier$
to
$\TestFunction\circ\IdentifyingIsometry$.
The change-of-variables identity for each fiber measure then implies
\[
 \begin{aligned}
 \int_\ComparisonCarrier\TestFunction\,
   \IntegrationDifferential(\IdentifyingIsometry\Pushforward\SelectedMeasure{\BaseCarrier}{\MetricSymbol})
 &=\int_{\MeasureFiber{\BaseCarrier}}
     \left(\int_\BaseCarrier\TestFunction\circ\IdentifyingIsometry\,
       \IntegrationDifferential\FiberMeasure{\BaseCarrier}{\FiberPoint}\right)
       \IntegrationDifferential\LawSection\MetricSpace{\BaseCarrier}{\MetricSymbol}(\FiberPoint)\\
 &=\int_{\MeasureFiber{\ComparisonCarrier}}
     \left(\int_\ComparisonCarrier\TestFunction\,
       \IntegrationDifferential\FiberMeasure{\ComparisonCarrier}{\FiberPoint}\right)
       \IntegrationDifferential\LawSection\MetricSpace{\ComparisonCarrier}{\ComparisonMetric}(\FiberPoint)\\
 &=\int_\ComparisonCarrier\TestFunction\,
       \IntegrationDifferential\SelectedMeasure{\ComparisonCarrier}{\ComparisonMetric}.
 \end{aligned}
\]
Uniqueness of the representing measure proves the required compatibility
with isometries.

\paraid{p-7ae020536d8d}
It remains to prove continuity as the underlying spaces vary.
We use the prescribed embeddings into
$\AmbientCarrier$
and identify each carrier with its image.
Fix a real continuous function
$\TestFunction\in\ContinuousFunctions(\AmbientCarrier)$.
We express the integral of
$\TestFunction$
against each selected measure using one bounded continuous extension
on a fixed space.

\textbf{Step 1.
A closed domain for the fiber integrals.}
\paraid{p-72ddde2e12f8}
Put
\[
 \SequenceParameterSpace=\{0\}\cup\{2^{-\SequenceIndex}\mid\SequenceIndex\in\NonnegativeIntegers\}.
\]
Set
\[
 \ParameterCarrier_0=\BaseCarrier.
\]
For each
$\SequenceIndex\in\NonnegativeIntegers$,
set
\[
 \ParameterCarrier_{2^{-\SequenceIndex}}=\BaseCarrier_\SequenceIndex.
\]
All carriers have the restricted ambient metric.
The map
$\ParameterTime\mapsto[\ParameterCarrier_\ParameterTime]$
from
$\SequenceParameterSpace$
to
$\GHSpace$
is continuous.
Here the brackets denote the isometry class with the restricted metric.
Consequently,
\[
 \FiberIntegralDomain
 =\{(\ParameterTime,\FiberPoint)\in\SequenceParameterSpace\times\InvariantMeasuredSpaces
       \mid\MeasureProjection(\FiberPoint)=[\ParameterCarrier_\ParameterTime]\}
\]
is closed,
since it is the equality set of two continuous maps to the metric space
$\GHSpace$.
Define a function
$\FiberIntegral\colon\FiberIntegralDomain\to\RealNumbers$
by
\[
 \FiberIntegral(\ParameterTime,\FiberPoint)
 =\int_{\ParameterCarrier_\ParameterTime}\TestFunction\,
       \IntegrationDifferential\FiberMeasure{\ParameterCarrier_\ParameterTime}{\FiberPoint}.
\]
Continuity at parameter
$0$
follows from \eqref{eq:fiber-pullbacks}.
At an isolated parameter,
continuity follows from \eqref{eq:fiber-pullbacks}
with the carrier fixed.
Moreover,
$|\FiberIntegral|\leq\norm{\TestFunction}_\infty$.

\textbf{Step 2.
Extension and weak convergence of the laws.}
\paraid{p-b2e0811e32ee}
The product
$\SequenceParameterSpace\times\InvariantMeasuredSpaces$
is metrizable and hence normal.
By the Tietze--Urysohn theorem (\Cref{thm:tietze}),
$\FiberIntegral$
extends to a continuous function
$\ExtendedFiberIntegral\colon\SequenceParameterSpace\times\InvariantMeasuredSpaces\to\RealNumbers$
with
$|\ExtendedFiberIntegral|\leq\norm{\TestFunction}_\infty$.
Write
\[
 \ParameterLaw_\ParameterTime=\LawSection([\ParameterCarrier_\ParameterTime]).
\]
Continuity of
$\LawSection$
implies
$\ParameterLaw_{2^{-\SequenceIndex}}\to\ParameterLaw_0$
weakly.
\Cref{lem:dirac-product-convergence},
with
$\CompactParameterSpace=\SequenceParameterSpace$
and
$\ProbabilityCarrier=\InvariantMeasuredSpaces$,
implies
\begin{equation}\label{eq:parameter-law-products}
 \DiracProbability_{2^{-\SequenceIndex}}\otimes\ParameterLaw_{2^{-\SequenceIndex}}
 \longrightarrow\DiracProbability_0\otimes\ParameterLaw_0
 \quad\text{weakly on }\SequenceParameterSpace\times\InvariantMeasuredSpaces.
\end{equation}

\textbf{Step 3.
Integrals of the averaged measures.}
\paraid{p-b288c7e8154f}
Each product measure in \eqref{eq:parameter-law-products} is concentrated on
$\FiberIntegralDomain$.
By \eqref{eq:average-functional} and \eqref{eq:measure-barycenter},
\[
 \begin{aligned}
 \int_\AmbientCarrier\TestFunction\,\IntegrationDifferential\SelectedMeasure{\BaseCarrier_\SequenceIndex}{\MetricSymbol_\SequenceIndex}
 &=\int\ExtendedFiberIntegral\,
   \IntegrationDifferential(\DiracProbability_{2^{-\SequenceIndex}}\otimes\ParameterLaw_{2^{-\SequenceIndex}})\\
 &\longrightarrow\int\ExtendedFiberIntegral\,
   \IntegrationDifferential(\DiracProbability_0\otimes\ParameterLaw_0)
 =\int_\AmbientCarrier\TestFunction\,\IntegrationDifferential\SelectedMeasure{\BaseCarrier}{\MetricSymbol}.
 \end{aligned}
\]
Since this holds for every real continuous
$\TestFunction$
on the compact ambient space
$\AmbientCarrier$,
the selected measures
$\SelectedMeasure{\BaseCarrier_\SequenceIndex}{\MetricSymbol_\SequenceIndex}$
converge weakly to
$\SelectedMeasure{\BaseCarrier}{\MetricSymbol}$
as asserted.
\end{proof}
